\par\noindent\textbf{(1)} 纯态密度算符 \(\rho_{AB}=|\Phi^+\rangle\langle\Phi^+|\) 展开为 \((|00\rangle\langle00|+|00\rangle\langle11|+|11\rangle\langle00|+|11\rangle\langle11|)/2\)。对子系统B求偏迹，交叉项消失，得到 \(\rho_A=I_2/2\)。
\par\noindent\textbf{(2)} 约化态本征值为 \(1/2,1/2\)，因此 \(S(\rho_A)=-\operatorname{Tr}\rho_A\ln\rho_A=\ln2\)，以比特为单位为1。局部纯度 \(\operatorname{Tr}\rho_A^2=1/2\)，而全态 \(\operatorname{Tr}\rho_{AB}^2=1\)，说明全局纯不意味着局部纯。
\par\noindent\textbf{(3)} 取 \(A_0=\sigma_z,A_1=\sigma_x\)，\(B_0=(\sigma_z+\sigma_x)/\sqrt2,B_1=(\sigma_z-\sigma_x)/\sqrt2\)。对于 \(|\Phi^+\rangle\)，\(\langle\sigma_z\otimes\sigma_z\rangle=\langle\sigma_x\otimes\sigma_x\rangle=1\)，交叉平均为0，所以 \(\langle A_0(B_0+B_1)+A_1(B_0-B_1)\rangle=2\sqrt2\)。
\par\noindent\textbf{(4)} 局域确定性模型的每次赋值给 \(|a_0(b_0+b_1)+a_1(b_0-b_1)|\le2\)，混合后仍不超过2；量子可达 \(2\sqrt2\)，但不能超过Tsirelson界。实际检测损耗和通信漏洞需要实验额外条件，否则观测值不能简单视为无条件排除所有隐变量模型。
